
\documentclass[12pt,a4paper]{article}
\usepackage[utf8]{inputenc}
\usepackage[T1]{fontenc}
\usepackage[brazil]{babel}
\usepackage{geometry}
\usepackage{graphicx}
\usepackage{fancyhdr}
\usepackage{booktabs}
\usepackage{xcolor}
\usepackage{tikz}
\usepackage{pgfplots}
\usepackage{hyperref}
\usepackage{amsmath}
\usepackage{enumitem}

\geometry{margin=2.5cm}
\pgfplotsset{compat=1.17}

% Configuração do cabeçalho
\pagestyle{fancy}
\fancyhf{}
\fancyhead[L]{\textbf{AEONCOSMA - Relatório de Segurança}}
\fancyhead[R]{\today}
\fancyfoot[C]{\thepage}

% Cores personalizadas
\definecolor{aeonblue}{RGB}{0,102,204}
\definecolor{aeongray}{RGB}{128,128,128}
\definecolor{critical}{RGB}{220,20,60}
\definecolor{warning}{RGB}{255,165,0}
\definecolor{success}{RGB}{34,139,34}

\title{
    \vspace{-2cm}
    \begin{center}
        \includegraphics[width=0.3\textwidth]{aeon_logo.png}\\[1cm]
        \textbf{\Huge AEONCOSMA}\\[0.5cm]
        \textbf{\Large Relatório de Segurança da Rede}\\[0.3cm]
        \textcolor{aeongray}{\large Sistema de Monitoramento Digital Twin}
    \end{center}
    \vspace{-1cm}
}

\date{02 de August de 2025}
\author{Sistema Automatizado AEONCOSMA}

\begin{document}

\maketitle
\thispagestyle{empty}

\newpage

\tableofcontents

\newpage

\section{Resumo Executivo}

Este relatório apresenta uma análise abrangente da segurança da rede AEONCOSMA, cobrindo o período de 26/07/2025 - 02/08/2025. O sistema de monitoramento detectou 0 padrões de segurança e analisou 0 nós da rede.

\subsection{Status Geral da Segurança}

\begin{center}
\begin{tikzpicture}
    \pie[text=legend, radius=2.5]{
        7/Crítico,
        5/Alto,
        10/Médio,
        20/Baixo
    }
\end{tikzpicture}
\end{center}

\textbf{Nível de Segurança Geral:} \textcolor{critical}{0.0\% - ATENÇÃO NECESSÁRIA}

\subsection{Principais Descobertas}

\begin{itemize}[leftmargin=*]
    \item Nenhuma ameaça crítica detectada
    \item Excelente disponibilidade da rede (>95\%)
\end{itemize}

\section{Análise de Ameaças Detectadas}

\subsection{Detecções por Categoria}

\begin{table}[h!]
\centering
\begin{tabular}{@{}lcc@{}}
\toprule
\textbf{Categoria de Ameaça} & \textbf{Detecções} & \textbf{Severidade Máxima} \\
\midrule
Ataques DDoS & 2 & \textcolor{aeonblue}{MÉDIO} \\
Falhas de Consenso & 1 & \textcolor{warning}{ALTO} \\
Anomalias de Rede & 3 & \textcolor{success}{BAIXO} \\
Tentativas de Intrusão & 0 & \textcolor{aeongray}{N/A} \\
\bottomrule
\end{tabular}
\caption{Resumo das ameaças detectadas por categoria}
\end{table}

\subsection{Timeline de Eventos Críticos}

\begin{center}
\begin{tikzpicture}
    \begin{axis}[
        width=\textwidth,
        height=8cm,
        xlabel={Tempo},
        ylabel={Severidade},
        grid=major,
        legend pos=north west
    ]
    CRITICAL_EVENTS_PLOT
    \end{axis}
\end{tikzpicture}
\end{center}

\section{Métricas de Performance da Rede}

\subsection{Latência da Rede}

\textbf{Latência Média:} 0.0 ms\\
\textbf{Latência Máxima:} 0.0 ms\\
\textbf{Percentil 95:} 0.0 ms

\begin{center}
\begin{tikzpicture}
    \begin{axis}[
        width=\textwidth,
        height=6cm,
        xlabel={Tempo},
        ylabel={Latência (ms)},
        grid=major
    ]
    LATENCY_PLOT
    \end{axis}
\end{tikzpicture}
\end{center}

\subsection{Utilização de Recursos}

\begin{table}[h!]
\centering
\begin{tabular}{@{}lccc@{}}
\toprule
\textbf{Recurso} & \textbf{Média (\%)} & \textbf{Máximo (\%)} & \textbf{Status} \\
\midrule
CPU & 0.0 & 0.0 & \textcolor{success}{BAIXO} \\
Memória & 0.0 & 0.0 & \textcolor{success}{BAIXO} \\
Rede & AVERAGE_NETWORK & MAX_NETWORK & \textcolor{success}{BAIXO} \\
\bottomrule
\end{tabular}
\caption{Utilização de recursos do sistema}
\end{table}

\section{Análise de Consenso}

\subsection{Participação no Consenso}

\textbf{Taxa de Participação Média:} 0.0\%\\
\textbf{Nós Ativos no Consenso:} 0\\
\textbf{Falhas de Consenso:} 0

\begin{center}
\begin{tikzpicture}
    \begin{axis}[
        width=\textwidth,
        height=6cm,
        xlabel={Tempo},
        ylabel={Taxa de Participação (\%)},
        grid=major,
        ymin=0,
        ymax=100
    ]
    CONSENSUS_PLOT
    \end{axis}
\end{tikzpicture}
\end{center}

\section{Detecção de Anomalias}

\subsection{Algoritmos de Detecção Utilizados}

\begin{enumerate}
    \item \textbf{Detecção de Padrões Simbólicos}: Análise de sequências de eventos
    \item \textbf{Análise Estatística}: Desvios das métricas baseline
    \item \textbf{Machine Learning}: Detecção de comportamentos anômalos
    \item \textbf{Análise de Entropia}: Medição da aleatoriedade dos eventos
\end{enumerate}

\subsection{Resultados da Detecção}


O sistema de detecção de anomalias analisou 0 padrões distintos e identificou 0 eventos anômalos.
A entropia atual da rede é 4.50, indicando um nível normal de aleatoriedade nos eventos.

\begin{itemize}
\item Padrões simbólicos detectados: 0
\item Eventos anômalos identificados: 0
\item Score de entropia: 4.50
\item Desvios estatísticos: 0
\end{itemize}


\section{Recomendações de Segurança}

\subsection{Ações Imediatas}

\item Verificar integridade dos nós validadores
\item Atualizar regras de firewall se necessário

\subsection{Melhorias de Longo Prazo}

\item Implementar algoritmos de machine learning para detecção proativa
\item Desenvolver sistema de resposta automática a incidentes
\item Estabelecer métricas de baseline mais robustas
\item Criar dashboard em tempo real para operadores
\item Implementar sistema de backup automatizado

\section{Conclusões}


O sistema apresenta boa segurança com algumas áreas para melhoria. O monitoramento contínuo e as análises automatizadas forneceram insights valiosos 
sobre o estado da rede. As recomendações apresentadas devem ser implementadas de acordo com 
suas prioridades para manter e melhorar a segurança operacional.

A implementação das medidas sugeridas resultará em maior resiliência e confiabilidade do 
sistema AEONCOSMA.


\section{Apêndices}

\subsection{Apêndice A: Dados Técnicos Detalhados}


\begin{verbatim}
Timestamp de Coleta: 2025-08-02T16:00:19.602958
Nós Totais: N/A
Nós Online: N/A
CPU Média: N/A%
Memória Média: N/A%
Latência Média: N/Ams
Taxa de Consenso: N/A
\end{verbatim}


\subsection{Apêndice B: Configurações do Sistema}


\begin{verbatim}
Sistema: AEONCOSMA Digital Twin Network
Versão: 1.0.0
Algoritmo de Consenso: Proof of Stake Modificado
Criptografia: SHA-256, ECDSA
Protocolos de Rede: TCP/IP, UDP, WebSocket
Monitoramento: Real-time com detecção de anomalias
\end{verbatim}


\end{document}
